\section{Background}

Local community service nonprofits provide essential services, such as food pantries, homeless shelters, substance use treatment, and at-risk youth programs. More than two-thirds of nonprofits report difficulty employing enough staff (Nonprofit Finance Fund 2025). Since January 2025, 73 percent of nonprofits have seen demand for their services increase, while 30 percent have reduced staff and 29 percent have reduced services (Grundhoefer, Smith Arrillaga, and Buteau 2026). Concern about staff burnout is widespread, with ninety percent of nonprofit leaders reporting concern about burnout due to funding challenges and rising demand for services (Martin 2026). In food, shelter, and relief organizations, that strain falls largely on caseworkers and social workers, who hold more than one in five jobs (U.S. Bureau of Labor Statistics 2026). In order to improve efficiency of routine tasks, nearly two-thirds of social workers already report using AI in their work (Borah et al. 2026).

Existing research about AI in the nonprofit community service sector primarily comes from adoption surveys (Smith Arrillaga, Grundhoefer, and Im 2025; NTEN and Bridgespan 2026; Shi and Azevedo 2026). Such surveys are largely aimed at the organizational level of AI usage, and do not translate usage to the level of individual employee time. In a 2026 survey of social workers, 64\% reported using AI in their job, mostly for documentation and reporting (Borah et al. 2026). Few social workers (\textless5\%) reported using AI for more time-consuming tasks such as case management -- ethical and data privacy concerns were cited as the most common reason (Borah et al. 2026). In short, while AI has the potential to offload burden at tasks such as reporting, clerical tasks, digital work, etc. that are part of community service occupations- time-sensitivity weighting highlights that such tasks take up a relatively small part of these occupations, and so AI exposure may be lower than indicated by task-weighting alone. Much of the nonprofit worker's time involves interpersonal interaction, which takes up more of the worker's time and is less AI-exposed.

In studying theoretical ability of AI to perform workforce tasks, Eloundou et al. (2024) rate each O*NET task on whether an LLM could cut its completion time by at least half at equal quality. A task is rated directly exposed (E1) if an LLM alone could achieve that reduction, partially exposed (E2) if the LLM could do it if additional software were built on the model, and not exposed (E0) otherwise. They developed exposure scores for every O*NET task in the database, providing a foundation for analyzing theoretical exposure to AI in the workforce. A problem with using these rankings directly when studying community service nonprofits, however, is that the Eloundou et al. scores are provided per task (i.e. ``driving'', ``submit reports'', ``assist in locating housing''), they do not take into account how much time workers actually spend per task.

In "Estimating Time Spent on Work Tasks," Hatgis-Kessell et al. propose a method for weighting tasks by time for every O*NET occupation, facilitating the computation of worker time exposed to AI, rather than just the share of worker tasks exposed. Hours are the product of O*NET-reported task frequency and a per-instance duration estimated from language-model pairwise comparisons, fitted to a 7-hour working day. I apply the Hatgis-Kessell et al. approach to get to the fraction of working time (rather than the fraction of tasks) to understand AI exposure for local community relief work.

Massenkoff and McCrory (2026) move from studying theoretical to observed exposure by measuring which O*NET tasks are currently seeing work-related use, measured from conversations with Claude. Again, information is provided for the entire O*NET database at the task level. Unlike surveys which rely on humans to accurately self-report, this data is drawn from recorded use. The research from Massenkoff and McCrory indicate that actual AI use in occupations falls well short of what AI could theoretically be doing. The authors attribute this gap to factors not captured by theoretical exposure studies, namely: legal constraints, the potential need for additional software, or the time it takes to check AI results.

Existing AI exposure studies each answer a different question. Theoretical exposure work from Eloundou et al. indicate which tasks AI might speed up. Time-weighting analysis from Hatgis-Kessell et al. looks at how much time each task takes. Observed exposure from Massenkoff and McCrory indicate where LLMs are already being used compared against where they theoretically could be used. This paper adds value to the current literature by combining three different economy-scale exposure studies and applying them deeply to understand the implications for local community food, housing, and emergency service work.
